\chapter*{Dedicatoria y Agradecimientos}
\label{cha:dedicatoria}

% Un tamaño de letra un poco menor y justificado normal se lee excelente en dos páginas
\begin{small} 
\itshape 
    \vspace{1em}
    % ==========================================
    % 1. AGRADECIMIENTOS CONJUNTOS
    % ==========================================
    \begin{center}
        \textbf{Agradecimientos Institucionales} \\[0.5em]
        \begin{minipage}{0.9\textwidth}
            \centering
            Nuestros sinceros agradecimientos a nuestro profesor guía, Nicolás Matus, por su apoyo, dirección y compromiso durante todo el proceso de esta tesis. También deseamos agradecer al profesor Cristian Rusu, profesor correferente, por sus constructivas observaciones y por la dedicación en la revisión de este trabajo de título, aportando enormemente a su desarrollo.

            Agradecemos a la Pontificia Universidad Católica de Valparaíso por brindarnos la oportunidad de formarnos académicamente en un ambiente de excelencia, y por el apoyo institucional que ha sido fundamental para alcanzar esta meta. Asimismo, extendemos nuestro agradecimiento a todos los profesores y personal administrativo que han contribuido a nuestra formación durante estos años, especialmente nombrando al profesor Doctor Ignacio Araya, por su dedicación y compromiso con la enseñanza y el desarrollo de sus estudiantes.
        \end{minipage}
    \end{center}
    
    \vspace{1em} % Espacio prudente para separar lo institucional de lo personal

    % ==========================================
    % 2. DEDICATORIA DE FABRIZZIO MURA
    % ==========================================
    \noindent\textbf{Fabrizzio Mura} \\
    \nopagebreak
    \indent Expreso mis más profundos agradecimientos a mi madre, que me impulsó a ser todo lo que fui, soy y siempre seré, la señora Bruna Helena Rosa Lavarello Paredes, que me ha guiado desde siempre y me ayuda de maneras que no se imagina a seguir adelante tanto en mi vida personal como profesional, además de mi perrita amada Luna.

    A la mujer que me impulsó a no atrasarme y seguir al día con mis estudios pese a mis dudas, me acompaña y apoya en los malos momentos y me hace ser mejor hombre y mejor persona cada día, Valeria Alexandra García Estella.

    A mis amigos, Diego Gallegos, Martín Gajardo, Nicolás Gómez, Matías Maturana, Ignacio Valdivia y Sebastián Valdivia por tantos años de amistad y compañía, que más que amigos se han vuelto una parte más de mi familia.

    A mis amigos en la universidad, que me han acompañado en esta formación académica desde distintos años pero que se han vuelto parte importante de mi vida, para hoy y siempre Joaquín Tapia, Matías Bugueño, Francisco Cortés y todos mis amigos y amigas cercanos que estuvieron ahí durante estos años tanto de estudio como de formación como persona, apreciándolos mucho como una etapa invaluable e inolvidable de mi vida.

    Y a todos los que he mencionado aquí o quizás me falte mencionar, por ser parte de mi vida, por reír, llorar, conversar, comprender y sobre todo permitirme formar parte de sus vidas, ya que cada persona aportó a formar una pequeña parte de mi ser y no sería lo mismo sin ellos(as), ya que no he llegado aquí solo, me han ayudado y he ayudado, me han cuidado y he cuidado, agradezco a cada persona con la que he interactuado por formar aunque sea por un segundo, una parte de mi vida.

    \vspace{1em} % Espacio de separación limpio antes del siguiente autor

    % ==========================================
    % 3. DEDICATORIA DE MATÍAS BUGUEÑO
    % ==========================================
    \noindent\textbf{Matías Bugueño} \\
    \nopagebreak
    \indent En primer lugar, deseo expresar mi más profundo y sincero agradecimiento a mi madre, Claudia Cecilia Bugueño Aravena, quien ha sido el pilar fundamental a lo largo de todo mi camino. Su apoyo incondicional y su fortaleza en los momentos más complejos de mi vida me permitieron llegar hasta este punto. Este logro es, sin duda alguna, tan suyo como mío.

    A mi hermanita perruna, Donatella, por su compañía y alegría constante, que me han brindado momentos de felicidad y distracción necesarios para mantener el equilibrio durante esta etapa.

    Asimismo, agradezco sinceramente a los profesores y compañeros de mi primer colegio, el Colegio Guillermo Zañartu Irigoyen, por cimentar las bases de mi formación académica y personal, impulsándome a dar el paso hacia una de las etapas más significativas de mi vida: mi ingreso al Instituto Nacional José Miguel Carrera. A esta última institución le agradezco enormemente la excelencia de su formación, sus valores y el rigor académico que me brindó. Las herramientas, conocimientos y oportunidades recibidas en sus aulas fueron determinantes para enfrentar con éxito la etapa universitaria. Honraré sus enseñanzas y llevaré siempre conmigo la convicción de su lema: \textit{Labor Omnia Vincit}. Como mención especial, al profesor José Garrido y a la profesora Mapi, quienes me brindaron todo su apoyo y fueron mi gran impulso en mi avance académico.

    A mis amigos, quienes me acompañaron de cerca en este recorrido, ofreciéndome su alegría, buena energía y un respaldo invaluable cuando más lo necesité; en especial a Antonia Luengo, Alejandra Hidalgo, Fabrizzio Mura, Sebastián Contreras, Joaquín Tapia y Thomas Molina.

    Finalmente, extiendo mi gratitud a mi familia por ser un soporte constante tanto en lo académico como en lo emocional. Su esfuerzo y cariño han sido el motor que me impulsó a cumplir esta meta, con una mención especial a Flavio Díaz, Julio Chelme, Julio Bugueño y María Bugueño por estar siempre presentes.

\end{small}
\addtocounter{page}{-1}